% ============================================================
%  CHAPTER 12 — Conclusion and Future Work
% ============================================================
\chapter{Conclusion and Future Work}
\label{ch:conclusion}

\section{Conclusion}

This thesis presented the design, implementation, and evaluation of a
\textbf{Vision--Language--Action agent} for the PAL Robotics \tiago{} mobile
manipulator.
The system enables the robot to accept natural-language commands — via voice or text
— and autonomously execute multi-step manipulation tasks by combining visual
perception, VLM-based reasoning, and a library of physical skills.

The main contributions of this work are:

\begin{enumerate}
  \item \textbf{A complete, deployable VLA agent} that integrates Google Gemini 2.0
        Flash, YOLO11n object detection, and nine physical skills into a coherent
        perceive--reason--act loop operating within the constraints of a Python 3.6
        / ROS Melodic Docker environment.

  \item \textbf{A robust 3-D object localisation pipeline} using RANSAC cylinder
        fitting on RGB-D point clouds with 5-frame averaging, achieving sub-5 mm
        positional stability for stationary objects.

  \item \textbf{A split-process YOLO microservice} that overcomes the Python 3.6
        constraint by deploying YOLO11n on the host machine and communicating via HTTP,
        achieving $\approx$18 ms inference with a clean separation between the
        AI inference layer and the ROS control layer.

  \item \textbf{An iterative grasping pipeline} that evolved through five versions,
        each addressing a specific failure mode, ultimately achieving a \todo{XX}\%
        grasp success rate across 27 trials at varied table positions.

  \item \textbf{A handover skill} that detects a person using YOLO and depth sensing,
        plans a collision-free arm trajectory toward them, and presents the grasped
        object at an ergonomic height, completing the full ``find--grasp--deliver''
        task chain.

  \item \textbf{A comprehensive safety system} that gates all skills via affordance
        checks, tracks consecutive failures, and enters a safe mode with automatic
        safe-pose recovery to prevent robot damage during failures.
\end{enumerate}

The system was evaluated on a real \tiago{} robot in a laboratory environment,
demonstrating end-to-end task completion of \todo{XX}\% for the full
``grab and hand over'' scenario, with a VLM skill selection accuracy of \todo{XX}\%.

\section{Future Work}

\subsection{Industrial-Style Trajectory Generation}

The current torso-descent grasp approach, while effective, is slow and produces
a kinematically suboptimal motion profile.
The planned v6 implementation will replace it with a MoveJ/MoveL trajectory generator
using Linear Segment with Parabolic Blends (LSPB) velocity profiles, providing smooth,
time-optimal motions that generalise to any object type.

\subsection{Navigation Integration}

Navigation is currently implemented but not integrated into the agent's skill loop.
Future work will close this gap by:
\begin{itemize}[noitemsep]
  \item Loading \texttt{NavigateToSkill} into the agent's skill library.
  \item Integrating location management (save/load named positions).
  \item Adding ``find the bottle on the shelf'' type tasks that require base movement.
\end{itemize}

\subsection{Object Permanence and Semantic Memory}

A scene graph updated across multiple time steps would allow the agent to remember
where objects were last seen and reason about object persistence.
This could be implemented as a Python dictionary (object name → last known 3-D
position) updated at each perception step.

\subsection{Local VLM Deployment}

Replacing the Gemini API with a locally deployed model (LLaVA 1.6, VILA, or InternVL)
would eliminate the internet dependency and reduce latency.
A GPU-equipped workstation connected to the robot via ROS would enable inference at
$<$200 ms per query.

\subsection{Adaptive Grasp Orientation}

Extending the grasping pipeline to estimate the object's principal axis and adapt
the approach orientation accordingly would enable grasping of non-cylindrical objects
lying in arbitrary orientations.
RANSAC-based axis detection on the depth point cloud is a natural extension of the
existing cylinder fitting approach.

\subsection{Open-Vocabulary Detection with GPU}

When GPU resources are available, replacing YOLO11n with YOLOE-26s-seg or
Grounding DINO would enable detection of arbitrary objects described by text,
removing the restriction to COCO-80 classes and enabling fully open-ended task
specifications such as ``pick up the red mug''.

\subsection{Ontology and Semantic Reasoning}

Integrating a lightweight semantic knowledge base (NetworkX scene graph or RDFLib
with SOMA ontology) would allow the VLM to query structured knowledge about object
affordances, spatial constraints, and task preconditions, moving beyond implicit
reasoning toward verifiable task-and-motion planning.

\section{Final Remarks}

The results of this work demonstrate that a zero-shot VLM-based agent, without any
robot-specific fine-tuning, can achieve reliable multi-step manipulation on a real
mobile manipulator.
The modular architecture — where perception, reasoning, and execution are cleanly
separated — makes it straightforward to improve individual components (faster detector,
better VLM, smoother trajectories) without restructuring the system.
This work lays a solid foundation for a more capable embodied AI agent as better
models and hardware become available.
